\section{Test Data}\label{sec:data}

\paragraph{Tasks and goals.} We use AgentDojo's workspace suite~\cite{agentdojo}: 40 user
tasks (benign traffic), such as ``Please send an email to each person mentioned in the TODO
list in the team meeting minutes file with the task they need to execute'', and 14 attacker
goals. An attack episode pairs a user task with a goal and one way of planting it
(Appendix Table~\ref{tab:attacks}): AgentDojo's four data attacks
(\texttt{important\_instructions}, \texttt{tool\_knowledge}, \texttt{direct},
\texttt{ignore\_previous}) and our static judge-aware variant \texttt{judge\_aware} plant text in emails, events
and files; our MCP attacks plant the goal in tool descriptions (\emph{tagged}: an Invariant
Labs-style \texttt{<IMPORTANT>} block~\cite{invariant2025}; \emph{plain}: one sentence) or in
parameter descriptions (tagged). AgentDojo's checkers then score each episode.
\emph{Attack success} (did the attacker's side effect happen?) is read from the final backend
state. \emph{Utility} (was the user's task done?) is read from the state and, for \UtilityReadsAnswer{} of the 40
tasks, also from the agent's answer.


\paragraph{Fairness.} Tasks whose number is divisible by 4 form a \emph{dev} split (10
tasks), the other 30 a \emph{test} split; all guard numbers in \S\ref{sec:results} are on test.
Each attack style uses the same 150 (task, goal) pairs (seed 7) for every configuration, so the
guard and no guard see identical cases. The final guard contains no pattern of any attack.
Every episode starts a fresh NLIP server and sends the task as an NLIP message to
\texttt{POST /nlip} through an in-process HTTP transport; the harness ran on our own laptop and
the models on Jetstream2. The harness disables the sign-in SE, so sign-in is not part of the
evaluated path. The 40 tasks, 14 goals and every planted text (560 cases per attack) are
exported as plain files. The poison is placed in the tools or data the reference solution
uses, so the model sees it whenever it follows that path, which favors the attacker.

\paragraph{Deviations from a clean hold-out.} Three things weaken the split. (i) Several
design choices came from two \emph{test} tasks, 13 and 25. The judge's wording and its 0.7
threshold were chosen from their judge inputs, into which five hand-written attacker calls were
inserted: the calls of goals 0, 1 and 2, a security-code email like goals 4--5, and sharing
file 13 with the attacker. The ranked evidence lines and the 1,500-character result preview
fixed false blocks on task 25, and span-only redaction and the 20-call cap came from
behaviour on task 13. (ii) Earlier
designs (the ablations) and two unreported partial pilot runs of the final design were run on
test before we chose it. (iii) The decision to drop the rules used the single-check
measurements, also on test data. \S\ref{sec:results} therefore also reports the headline
numbers without tasks 13 and 25.
